% Bagian: sets-functions-relations
% Bab: size-of-sets
% Subbagian: reduction-alt

\documentclass[../../../include/open-logic-section]{subfiles}

\begin{document}

\olfileid[id]{sfr}{siz}{red-alt}

\olsection{Reduksi}

\begin{editorial}
  Bagian ini membuktikan ketakterhitungan dengan reduksi, selaras dengan
  hasil-hasil dalam \olref[nen-alt]{sec}. Versi alternatif yang sedikit lebih
  terperinci dan selaras dengan hasil-hasil dalam \olref[nen]{sec} diberikan
  dalam \olref[red]{sec}.
\end{editorial}

Kita telah membuktikan bahwa $\Bin^\omega$ !!{nonenumerable} dengan argumen
diagonalisasi. Kita menggunakan argumen diagonalisasi serupa untuk menunjukkan
bahwa $\Pow{\Nat}$ !!{nonenumerable}. Namun, berikut cara lain untuk
membuktikan bahwa $\Pow{\Nat}$ !!{nonenumerable}: tunjukkan bahwa \emph{jika
$\Pow{\Nat}$ !!{enumerable}, maka $\Bin^\omega$ juga !!{enumerable}}. Karena
kita tahu bahwa $\Bin^\omega$ !!{nonenumerable}, akan mengikuti bahwa
$\Pow{\Nat}$ juga demikian.

Ini disebut \emph{mereduksi} satu masalah ke masalah lain. Dalam kasus ini,
kita mereduksi masalah mengenumerasi $\Bin^\omega$ ke masalah mengenumerasi
$\Pow{\Nat}$. Penyelesaian bagi masalah yang kedua---suatu enumerasi
$\Pow{\Nat}$---akan menghasilkan penyelesaian bagi masalah yang
pertama---suatu enumerasi $\Bin^\omega$.

Untuk mereduksi masalah mengenumerasi suatu himpunan~$B$ ke masalah
mengenumerasi suatu himpunan~$A$, kita memberikan cara untuk mengubah suatu
enumerasi~$A$ menjadi enumerasi~$B$. Cara termudah ialah mendefinisikan
!!a{surjection} $f\colon A \to B$. Jika $x_1$, $x_2$, \dots{} mengenumerasi
~$A$, maka $f(x_1)$, $f(x_2)$, \dots{} akan mengenumerasi~$B$. Dalam kasus
kita, kita mencari !!a{surjection} $f\colon \Pow{\Nat} \to \Bin^\omega$.

\begin{prob}
Tunjukkan bahwa jika terdapat fungsi !!{injective} $g\colon B \to A$ dan
$B$~adalah !!{nonenumerable}, maka $A$ juga demikian. Lakukan ini
dengan menunjukkan cara menggunakan~$g$ untuk mengubah suatu enumerasi~$A$
menjadi enumerasi~$B$.
\end{prob}

\begin{proof}[Bukti reduksi untuk {\olref[nen-alt]{thm:nonenum-pownat}}]
Untuk melakukan reduksi, andaikan $\Pow{\Nat}$ adalah !!{enumerable}, sehingga
terdapat suatu enumerasi darinya, yaitu $N_{1}$, $N_{2}$, $N_{3}$, \dots

Definisikan fungsi $f \colon \Pow{\Nat} \to \Bin^\omega$ dengan menetapkan
$f(N)$ sebagai untai $s$ sedemikian sehingga $s(n) = 1$ jika dan hanya jika
$n \in N$, dan $s(n) = 0$ jika tidak.

Ini jelas mendefinisikan suatu fungsi, karena jika $N \subseteq \Nat$, maka
setiap $n \in \Nat$ merupakan !!a{element} dari $N$ atau bukan. Sebagai
contoh, himpunan bilangan asli genap $2\Nat = \Setabs{2n}{n \in \Nat} =
\{0,2, 4, 6, \dots\}$ dipetakan ke untai $1010101\dots$; $\emptyset$ dipetakan
ke $0000\dots$; dan $\Nat$ dipetakan ke $1111\dots$.

Fungsi itu juga !!{surjective}: setiap untai $0$ dan $1$ bersesuaian dengan
suatu himpunan bilangan asli, yakni himpunan yang beranggotakan bilangan asli
yang bersesuaian dengan posisi tempat untai tersebut memuat~$1$. Secara lebih
saksama, jika $s \in \Bin^\omega$, definisikan $N \subseteq \Nat$ sebagai:
\[
N = \Setabs{n \in \Nat}{s(n) = 1}
\]
Maka $f(N) = s$, sebagaimana dapat diperiksa dari definisi~$f$.

Sekarang perhatikan daftar
\[
f(N_1), f(N_2), f(N_3), \dots
\]
Karena $f$ !!{surjective}, setiap anggota $\Bin^\omega$ harus muncul sebagai
nilai~$f$ untuk suatu argumen, sehingga harus muncul dalam daftar. Jadi,
daftar ini harus mengenumerasi seluruh~$\Bin^\omega$.

Dengan demikian, jika $\Pow{\Nat}$ !!{enumerable}, maka $\Bin^\omega$ juga
!!{enumerable}. Namun, $\Bin^\omega$ adalah !!{nonenumerable}
(\olref[nen-alt]{thm:nonenum-bin-omega}). Jadi, $\Pow{\Nat}$ adalah
!!{nonenumerable}.
\end{proof}

%\begin{explain}
%Kita mudah keliru mengenai arah suatu reduksi.
%Sebagai contoh, fungsi !!a{surjective} $g \colon \Bin^\omega \to X$
%\emph{tidak} membuktikan bahwa $X$ !!{nonenumerable}. (Perhatikan $g
%\colon \Bin^\omega \to \Bin$ yang didefinisikan oleh $g(s) = s(1)$, yaitu
%fungsi yang memetakan untai $0$ dan $1$ ke !!{element} pertamanya. Fungsi itu
%surjektif karena beberapa untai dimulai dengan $0$ dan beberapa lainnya
%dengan $1$. Namun, $\Bin$ hingga.) Perhatikan pula bahwa fungsi $f$ harus
%surjektif; jika tidak, argumennya tidak berlaku:
%$f(x_1)$, $f(x_2)$, \dots{} tidak dijamin memuat semua !!{element}s dari~$Y$.
%Sebagai contoh, $h\colon \Nat \to
%\Bin^\omega$ yang didefinisikan oleh
%\[
%h(n) = \underbrace{000\dots0}_{\text{$n$ buah $0$}}
%\]
%adalah suatu fungsi, tetapi $\Nat$ adalah !!{enumerable}.
%\end{explain}

\begin{prob}\label{sfr:siz:red-alt:prob:nat-nat}
  Tunjukkan dengan argumen reduksi bahwa himpunan~$X$ dari semua fungsi
  $f\colon \Nat \to \Nat$ adalah !!{nonenumerable} (Petunjuk: berikan fungsi
  surjektif dari $X$ ke~$\Bin^\omega$.)
\end{prob}

\begin{prob}
Tunjukkan dengan argumen reduksi bahwa himpunan dari semua \emph{himpunan}
pasangan bilangan asli, yakni $\Pow{\Nat \times \Nat}$, adalah
!!{nonenumerable}.
\end{prob}

\begin{prob}
Tunjukkan dengan argumen reduksi bahwa $\Nat^\omega$, himpunan barisan tak
hingga bilangan asli, adalah !!{nonenumerable}.
\end{prob}

%\begin{prob}
%Misalkan $P$ adalah himpunan fungsi dari $\Nat$ ke himpunan $\{0\}$, dan misalkan $Q$ adalah himpunan fungsi \emph{parsial}
%dari himpunan bilangan bulat positif ke himpunan $\{0\}$. Tunjukkan
%bahwa $P$~adalah !!{enumerable} dan $Q$~tidak. (Petunjuk: reduksi masalah
%mengenumerasi $\Bin^\omega$ ke masalah mengenumerasi~$Q$.)
%\end{prob}

\begin{prob}
Misalkan $S$ adalah himpunan semua !!{surjection}s dari $\Nat$ ke himpunan
$\{0,1\}$, yakni $S$ terdiri atas semua !!{surjection}s~$f \colon \Nat \to
\Bin$. Tunjukkan bahwa $S$ adalah !!{nonenumerable}.
\end{prob}

\begin{prob}
Tunjukkan bahwa himpunan~$\Real$ dari semua bilangan real adalah
!!{nonenumerable}.
\end{prob}

\end{document}
